\section*{Chapitre \ref{ch:b2:corrosion} --- Corrosion et protection}

\begin{solution}{exo:b2:corrosion:1}
$j = \qty{0.020}{A/m^2}$~: $de/dt = 0.020 \times 65.38/(2 \times 96\,485 \times 7.13
\times 10^6) = \qty{9.5e-13}{m/s}$, soit \qty{0.030}{mm} par an.
\end{solution}

\begin{solution}{exo:b2:corrosion:2}
Boulons en acier sur du cuivre~: le fer, moins noble, est l'anode, et les petits boulons
sur une grande cathode de cuivre se corrodent vite. Clous en zinc dans de l'acier~: le zinc se corrode
et protège l'acier qui l'entoure.
\end{solution}

\begin{solution}{exo:b2:corrosion:3}
Dans l'étroit interstice entre les plaques et sous la tête du boulon, où l'eau
se renouvelle lentement et où son dioxygène est vite épuisé~: par \omterm{def:b2:corrosion:differential}{aération différentielle},
ces zones confinées deviennent des anodes, les surfaces dégagées des cathodes.
\end{solution}

\begin{solution}{exo:b2:corrosion:4}
La boîte en fer-blanc~: l'étain, plus noble que le fer, fait de l'acier mis à nu l'anode d'une
grande cathode. Sur le seau galvanisé, c'est le zinc qui se corrode à la place de
l'acier.
\end{solution}

\begin{solution}{exo:b2:corrosion:5}
L'agitation amincit la \omterm{def:b2:current-potential-curves:limiting-current}{couche de diffusion}~: le palier du dioxygène s'élève, la courbe anodique
du fer le rencontre plus haut, et le potentiel comme le \omterm{def:b2:corrosion:uniform}{courant de corrosion}
augmentent. Sans dioxygène, la seule réaction cathodique qui subsiste dans une eau
neutre est la réduction lente de l'eau~: le \omterm{def:b2:corrosion:uniform}{potentiel de corrosion} baisse et le
\omterm{def:b2:corrosion:uniform}{courant de corrosion} devient très faible.
\end{solution}

\begin{solution}{exo:b2:corrosion:6}
$t = 0.90 \times 5000 \times 2 \times 96\,485/(65.38 \times 0.10) = \qty{1.33e8}{s}$,
environ 4.2~ans.
\end{solution}

\begin{solution}{exo:b2:corrosion:7}
Charge par kilogramme, $nF/M$~: zinc \qty{820}{A.h/kg}, magnésium \qty{2.21e3}{A.h/kg},
aluminium \qty{2.98e3}{A.h/kg}. Potentiels standard~: \ce{Zn^2+}/\ce{Zn} $-0.76$,
\ce{Mg^2+}/\ce{Mg} $-2.36$, \ce{Al^3+}/\ce{Al} \qty{-1.68}{V}. Dans un sol sec et résistant,
le courant doit traverser une grande résistance~: le magnésium, bien plus
bas que le fer, offre la plus grande tension motrice.
\end{solution}

\begin{solution}{exo:b2:corrosion:8}
$U = 2.0 \times 4.0 + 1.5 = \qty{9.5}{V}$~; puissance \qty{19}{W}~; par an, $19 \times
8766 = \qty{1.7e5}{W.h}$, \qty{1.7e2}{kW.h}.
\end{solution}

\begin{solution}{exo:b2:corrosion:9}
À la ligne de flottaison, le métal est mouillé par une eau riche en dioxygène~: c'est une cathode~; juste
au-dessous, l'eau est moins aérée et la zone devient l'anode de la pile
formée avec la ligne de flottaison~: l'acier s'y corrode, selon une bande.
\end{solution}

\begin{solution}{exo:b2:corrosion:10}
Dans l'eau aérée, le \omterm{def:b2:corrosion:uniform}{potentiel de corrosion} de l'acier inoxydable tombe sur son
palier de passivité~: un courant minuscule, le film se reformant de lui-même. Les ions chlorure rompent
le film en ses points faibles~; le métal mis à nu y est une petite anode face à une
immense cathode passive, si bien que sa densité de courant est énorme. Dans la piqûre, les
ions métalliques s'hydrolysent et acidifient la solution, des ions chlorure y sont attirés pour équilibrer
la charge, et le film ne peut pas se reformer~: la piqûre se creuse.
\end{solution}

\begin{solution}{exo:b2:corrosion:11}
Le dioxygène est réduit sur la totalité des \qty{22}{cm^2}~: \omterm{def:b2:current-potential-curves:ie-curve}{courant cathodique} $22 \times 5 =
\qty{110}{\micro A}$, entièrement fourni par l'oxydation des \qty{2}{cm^2} d'acier
proches du raccord, $j = \qty{55}{\micro A/cm^2}$. La perte vaut $5.5$ fois celle du fer
à \qty{10}{\micro A/cm^2}, soit environ \qty{0.64}{mm} par an.
\end{solution}

\begin{solution}{exo:b2:corrosion:12}
Le fer est immunisé au-dessous de la droite \ce{Fe^2+}/\ce{Fe}, $-0.41 + 0.0296\log 10^{-6} =
\qty{-0.59}{V}$~: maintenu au-dessous d'environ \qty{-0.6}{V}, l'acier ne se dissout plus. Beaucoup
plus bas, l'eau est réduite sur l'acier à une vitesse croissante (la droite \ce{H+}/\ce{H2} est à
\qty{-0.41}{V} à pH 7 et la \omterm{def:b2:current-potential-curves:overpotential}{surtension} sur l'acier vaut quelques dixièmes de volt)~:
du courant est gaspillé et le dihydrogène décolle le revêtement.
\end{solution}

\begin{solution}{pb:b2:corrosion:1}
\textbf{1.} \ce{Fe -> Fe^2+ + 2e-}~; \ce{O2 + 2H2O + 4e- -> 4OH-}.
\textbf{2.} La réduction du dioxygène, lente et parvenue à son palier de diffusion,
limite le \omterm{def:b2:current-potential-curves:ie-curve}{courant cathodique}~; la courbe anodique du fer la rencontre sur ce palier.
\textbf{3.} $j = \qty{0.10}{A/m^2}$~: $0.10 \times 3.156 \times 10^7/(2 \times 96\,485)
\times 55.845 = \qty{913}{g}$ par mètre carré et par an.
\textbf{4.} $913/7.87 \times 10^6 = \qty{1.16e-4}{m}$, \qty{0.12}{mm} par an.
\textbf{5.} $4/0.116 = 34$~ans.
\textbf{6.} La corrosion n'est pas uniforme~: interstices, différences d'aération
entre les sols et défauts du revêtement concentrent le \omterm{def:b2:current-potential-curves:ie-curve}{courant anodique} sur
de petites surfaces, où des piqûres percent la paroi bien avant.
\textbf{7.} $E^\circ = \Delta_f G^\circ/(nF)$ pour chaque couple~: \ce{Fe} $-0.41$, \ce{Zn} $-0.76$,
\ce{Mg} $-2.36$, \ce{Al} \qty{-1.68}{V}.
\textbf{8.} Le zinc, le magnésium et l'aluminium, tous situés au-dessous du fer.
\textbf{9.} L'aluminium se passive~: son film d'oxyde arrête sa dissolution et il ne
débite alors que peu de courant, sauf s'il est allié pour empêcher la formation du film (ce qui fonctionne
dans l'eau de mer, pas dans les sols).
\textbf{10.} Le courant doit traverser la résistance du sol~: la tension
motrice entre l'anode et l'acier protégé vaut environ \qty{1.9}{V} pour le
magnésium contre \qty{0.35}{V} pour le zinc (d'après les potentiels standard).
\textbf{11.} $3.156 \times 10^7 \times 24.305/(2 \times 96\,485 \times 0.50) =
\qty{7.95}{kg}$ par ampère et par an.
\textbf{12.} $\pi \times 0.30 \times 10\,000 = \qty{9.42e3}{m^2}$.
\textbf{13.} \qty{94.2}{m^2}.
\textbf{14.} $0.020 \times 94.2 = \qty{1.88}{A}$.
\textbf{15.} $1.885 \times 20 \times 3.156 \times 10^7 = \qty{1.19e9}{C}$.
\textbf{16.} $1.19 \times 10^9 \times 24.305/(2 \times 96\,485 \times 0.50) =
\qty{3.00e5}{g}$, \qty{300}{kg} (ou $1.885 \times 20 \times 7.95$).
\textbf{17.} $300/15 = 20$ anodes.
\textbf{18.} La résistance du sol limite la distance sur laquelle une seule anode peut pousser
son courant le long de la conduite~: des anodes réparties la protègent uniformément.
\textbf{19.} Un redresseur fait circuler un courant depuis une anode inerte enterrée dans un lit
d'anode, à travers le sol, jusqu'à la conduite reliée à son pôle négatif.
\textbf{20.} $1.885 \times 5.0 + 1.5 = \qty{10.9}{V}$.
\textbf{21.} $10.9 \times 1.885 = \qty{20.6}{W}$~; $20.6 \times 8766 = \qty{1.8e5}{W.h}$,
\qty{180}{kW.h} par an.
\textbf{22.} L'eau est réduite sur la conduite~: du dihydrogène se forme, décolle le revêtement et
peut fragiliser l'acier, et du courant est gaspillé.
\textbf{23.} $E < -0.41 + 0.0296\log(10^{-6}) = \qty{-0.59}{V}$.
\textbf{24.} Vingt ans de protection exigent $\boldsymbol{\approx \qty{3.0e2}{kg}}$
d'anodes de magnésium.
\end{solution}
